% Method
\section{Method: Weighted Multi-Token Prediction (WMTP)}

\subsection{Problem Setup and Notation}

Let $x$ denote the input (problem description) and $y_{1:T}$ the output sequence (solution code). We train a policy $\pi_\theta$ with MTP heads $k \in \{1, \ldots, K\}$, where each head produces logits $z_{t,k}$ for predicting token $y_{t+k}$.

Following the Alpaca template, we apply \textbf{loss masking}: instruction and input portions have labels set to $-100$ (ignored in loss computation), training only on the output (solution) tokens.

\subsection{Baseline: Multi-Token Prediction Objective}

The standard MTP objective averages cross-entropy loss across all heads and tokens:
\begin{equation}
    \mathcal{L}_{\text{MTP}} = \frac{1}{K} \sum_{k=1}^{K} \sum_t \text{CE}(z_{t,k}, y_{t+k})
    \label{eq:mtp_baseline}
\end{equation}

This treats all tokens equally, regardless of their importance for solution correctness.

\subsection{Phase 1: Independent Critic via Pairwise Ranking}

\subsubsection{Data Construction}

We construct training pairs from the CodeContests dataset:
\begin{itemize}[nosep]
    \item For each problem, pair correct solutions with incorrect solutions
    \item Apply \textbf{length-balanced stratified sampling} using 17 bins to prevent length bias
    \item Limit to \texttt{max\_pairs\_per\_problem=60} to ensure diversity
    \item Total: 200K pairs (400K samples)
\end{itemize}

\subsubsection{Critic Architecture}

The critic model consists of:
\begin{itemize}[nosep]
    \item \textbf{Backbone}: Sheared-LLaMA 2.7B (HuggingFace LlamaModel)
    \item \textbf{Value head}: MLP with hidden dimension $\times 4 \rightarrow 1$
    \item \textbf{Output}: Token-level value $V_t \in \mathbb{R}$ for each position
    \item \textbf{Fine-tuning}: LoRA with rank=32, $\alpha$=64, dropout=0.1
\end{itemize}

\subsubsection{Pairwise Ranking Loss (Bradley-Terry)}

Rather than predicting absolute values via MSE regression, we use pairwise ranking to learn relative preferences. For a correct solution ($+$) and incorrect solution ($-$), we compute sequence-level values by averaging over output tokens:
\begin{equation}
    \bar{V}^{+} = \frac{1}{|\mathcal{O}^{+}|} \sum_{t \in \mathcal{O}^{+}} V_t^{+}, \quad
    \bar{V}^{-} = \frac{1}{|\mathcal{O}^{-}|} \sum_{t \in \mathcal{O}^{-}} V_t^{-}
\end{equation}

where $\mathcal{O}$ denotes the set of output token positions (excluding instruction/input).

The Bradley-Terry pairwise ranking loss encourages the correct solution's value to exceed the incorrect solution's:
\begin{equation}
    \mathcal{L}_{\text{pair}} = -\log \sigma(\bar{V}^{+} - \bar{V}^{-})
    \label{eq:pairwise_loss}
\end{equation}

This formulation is more appropriate for code where \textit{relative} correctness is easier to judge than absolute quality scores.

% Pairwise ranking diagram placeholder
\begin{figure}[h]
    \centering
    \fbox{\parbox{0.85\textwidth}{
        \centering
        \textbf{[Figure 2: Pairwise Ranking Training]}\\[1em]
        Correct solution $\xrightarrow{\text{Critic}}$ $V_t^{+}$ $\xrightarrow{\text{avg}}$ $\bar{V}^{+}$\\[0.3em]
        Incorrect solution $\xrightarrow{\text{Critic}}$ $V_t^{-}$ $\xrightarrow{\text{avg}}$ $\bar{V}^{-}$\\[0.5em]
        $\mathcal{L}_{\text{pair}} = -\log \sigma(\bar{V}^{+} - \bar{V}^{-})$
    }}
    \caption{Pairwise ranking training for the critic. The model learns that correct solutions should have higher average token values than incorrect solutions for the same problem.}
    \label{fig:pairwise}
\end{figure}

\subsection{Phase 2: Token Advantage Estimation with GAE}

\subsubsection{Objective}

Given the trained critic, we estimate token-level importance (credit assignment) from value changes across the sequence. Tokens that increase the value estimate contribute positively to correctness; tokens that decrease it contribute negatively.

\subsubsection{Generalized Advantage Estimation}

The standard GAE formulation \citep{schulman2015high} defines:
\begin{align}
    \delta_t &= r_t + \gamma V(s_{t+1}) - V(s_t) \label{eq:td_error} \\
    A_t &= \sum_{l \geq 0} (\gamma \lambda)^l \delta_{t+l} \label{eq:gae}
\end{align}

\subsubsection{Simplification for Code Generation}

Code generation is an \textbf{episodic task} with sparse rewards. We make the following simplifications:

\begin{itemize}[nosep]
    \item \textbf{No intermediate rewards}: $r_t = 0$ for all $t < T$
    \item \textbf{No terminal reward signal}: The value function itself encodes correctness likelihood (trained via pairwise ranking), so we do not add explicit $r_T$
    \item \textbf{No discounting}: $\gamma = 1.0$ (episodic task, no preference for earlier rewards)
\end{itemize}

This yields the simplified TD-error:
\begin{equation}
    \delta_t = \gamma V(s_{t+1}) - V(s_t) = V_{t+1} - V_t
    \label{eq:simplified_td}
\end{equation}

which represents the \textbf{pure value change} caused by token $t$.

\textbf{Hyperparameter}: We use $\lambda = 0.95$ (GAE standard recommendation), providing a bias-variance tradeoff between TD(0) ($\lambda=0$) and Monte Carlo ($\lambda=1$).

\subsubsection{Interpretation}

\begin{itemize}[nosep]
    \item $V$ \textbf{increases} $\Rightarrow$ Token contributes to higher correctness probability
    \item $V$ \textbf{decreases} $\Rightarrow$ Token contributes to lower correctness probability
\end{itemize}

\subsection{AWR-Style Weighting with Stabilization}

\subsubsection{Whitening Normalization}

Raw advantages can have unbounded scale, causing training instability. We apply whitening:
\begin{equation}
    \hat{A}_t = \frac{A_t - \mu}{\sigma + \epsilon}
    \label{eq:whitening}
\end{equation}

where $\mu$ and $\sigma$ are computed using \textbf{EMA running statistics} across batches:
\begin{align}
    \mu_{\text{new}} &= (1 - \alpha) \cdot \mu_{\text{old}} + \alpha \cdot \mu_{\text{batch}} \\
    \sigma_{\text{new}} &= (1 - \alpha) \cdot \sigma_{\text{old}} + \alpha \cdot \sigma_{\text{batch}}
\end{align}

with momentum $\alpha = 0.1$ and warmup period of 10 steps.

\subsubsection{Exponential Weighting with Clipping}

Following AWR, we convert normalized advantages to exponential weights:
\begin{equation}
    w_t = \exp\left(\frac{\hat{A}_t}{\beta}\right)
    \label{eq:exp_weight}
\end{equation}

and apply clipping to prevent extreme values:
\begin{equation}
    w_t \leftarrow \text{clip}(w_t, w_{\min}, w_{\max})
    \label{eq:clipping}
\end{equation}

\textbf{Hyperparameters}: $\beta = 1.0$, $[w_{\min}, w_{\max}] = [0.1, 3.0]$

\subsubsection{Distributed Training Synchronization}

For multi-GPU training with FSDP, we synchronize EMA statistics across ranks using \texttt{all\_reduce} to ensure consistent normalization.

\subsection{Final WMTP Objective}

The weighted MTP loss modulates each token's contribution by its importance weight:
\begin{equation}
    \boxed{
    \mathcal{L}_{\text{WMTP}} = \frac{1}{K} \sum_{k=1}^{K} \sum_t w_t \cdot \text{CE}(z_{t,k}, y_{t+k})
    }
    \label{eq:wmtp}
\end{equation}

Tokens with high advantage (contributing to correctness) receive higher weights; tokens with low or negative advantage receive lower weights.

\subsection{Implementation Details}

\subsubsection{Data Pipeline}

\begin{itemize}[nosep]
    \item \textbf{Template}: Alpaca format with instruction/input/output sections
    \item \textbf{Loss masking}: \texttt{labels=-100} for instruction and input tokens
    \item \textbf{Maximum length}: 2048 tokens
\end{itemize}

\subsubsection{Training Infrastructure}

\begin{itemize}[nosep]
    \item \textbf{Hardware}: NVIDIA H100 80GB $\times$ 4
    \item \textbf{Distributed training}: FSDP with \texttt{FULL\_SHARD} (ZeRO-3)
    \item \textbf{Precision}: BF16 mixed precision
    \item \textbf{Memory optimization}: Activation checkpointing
    \item \textbf{Policy LoRA}: rank=64, $\alpha$=128, dropout=0.05
\end{itemize}

\subsubsection{Inference Cost}

\begin{itemize}[nosep]
    \item \textbf{Training}: Phase 2 requires critic forward pass (approximately 1.3$\times$ training cost)
    \item \textbf{Inference}: \textbf{Zero additional cost}---the critic is not used at test time
\end{itemize}

\subsection{Algorithm Summary}

\begin{algorithm}[h]
\caption{Weighted Multi-Token Prediction (WMTP)}
\label{alg:wmtp}
\begin{algorithmic}[1]
\Require Critic $V_\phi$ (trained in Phase 1), Policy $\pi_\theta$, hyperparameters $\gamma, \lambda, \beta, w_{\min}, w_{\max}$
\For{each training batch}
    \State Tokenize input and apply loss mask
    \State Compute token values $V_t = V_\phi(x, y_{1:t})$ \Comment{Frozen critic}
    \State Compute TD-errors $\delta_t = \gamma V_{t+1} - V_t$
    \State Compute GAE advantages $A_t = \sum_{l \geq 0} (\gamma\lambda)^l \delta_{t+l}$
    \State Update EMA: $\mu, \sigma \leftarrow \text{EMA}(\mu, \sigma, A)$
    \State Normalize: $\hat{A}_t = (A_t - \mu) / (\sigma + \epsilon)$
    \State Compute weights: $w_t = \text{clip}(\exp(\hat{A}_t / \beta), w_{\min}, w_{\max})$
    \State Compute MTP logits $z_{t,k}$ for all heads $k$
    \State Compute weighted loss: $\mathcal{L} = \frac{1}{K} \sum_k \sum_t w_t \cdot \text{CE}(z_{t,k}, y_{t+k})$
    \State Update $\theta$ via gradient descent on $\mathcal{L}$
\EndFor
\end{algorithmic}
\end{algorithm}
